\begin{abstract}
I \textit{Large Language Models} (LLM) basano il loro funzionamento sull'elaborazione di rappresentazioni vettoriali ad alta dimensionalità, rendendo il loro processo decisionale difficile da interpretare. Il presente lavoro propone un approccio esplorativo volto a tradurre queste dinamiche interne in una rappresentazione visiva e geometrica tangibile. Attraverso l'implementazione di un ambiente per l'esplorazione tridimensionale dello spazio latente dei \textit{last hidden states}, il progetto si pone l'obiettivo di misurare e rendere direttamente osservabile l'impatto spaziale di alcune tecniche di \textit{prompting} avanzato.

L'analisi spaziale delle traiettorie di generazione ha mostrato come il contesto fornito nel prompt modifichi la distribuzione dei token, inducendo la formazione di \textit{cluster} semantici distinti. I risultati confermano che istruzioni ben strutturate non si limitano a orientare l'output testuale, ma guidano il modello verso specifiche regioni dello spazio latente, riducendo l'entropia nella scelta dei token successivi. Il lavoro fornisce uno strumento pratico per esplorare la ``scatola nera'' della decodifica probabilistica attraverso traiettorie spaziali misurabili, ponendo le basi per future esplorazioni di architetture \textit{Retrieval-Augmented Generation} (RAG) e di pattern di ragionamento più complessi.
\end{abstract}